\subsubsection*{\texttt{27-index-hierarchical-markdown}}
\label{app:samples:27_index_hierarchical_markdown}

\paragraph{Q3: How long does the Haste spell last?}
\noindent \textbf{Corr.:}~5.00/5\quad \textbf{Faith.:}~1.00\quad \textbf{Relev.:}~1.00\quad \textbf{Ctx.Rel.:}~1.00\\

\noindent\textbf{Antwort:} The Haste spell lasts for the duration of concentration, up to 1 minute. (From the provided rules.)

\noindent\textbf{Abgerufene Chunks:}
\begin{itemize}[noitemsep, topsep=2pt]
  \item[\textbf{\#1}] \texttt{[0.0167]} \textit{Spells/Haste} -- \#\#\# Haste  *3rd-level transmutation*  **Casting Time:** 1 action  **Range:** 30 \ldots
  \item[\textbf{\#2}] \texttt{[0.0164]} \textit{Spells/Haste} -- \#\#\# Haste  *3rd-level transmutation*  **Casting Time:** 1 action  **Range:** 30 \ldots
  \item[\textbf{\#3}] \texttt{[0.0161]} \textit{Spells/Modify Memory} -- If it takes any damage or is targeted by another spell, this spell ends, and non\ldots
  \item[\textbf{\#4}] \texttt{[0.0159]} \textit{Spells/Wish} -- After enduring that stress, each time you cast a spell until you finish a long r\ldots
  \item[\textbf{\#5}] \texttt{[0.0156]} \textit{Spells/Bestow Curse} -- If you use a spell slot of 7th level or higher, the duration is 24 hours. If you\ldots
\end{itemize}
\vspace{4pt}
\hrule
\vspace{6pt}

\paragraph{Q5: Can a grappled creature cast spells?}
\noindent \textbf{Corr.:}~3.00/5\quad \textbf{Faith.:}~0.00\quad \textbf{Relev.:}~0.00\quad \textbf{Ctx.Rel.:}~—\\

\noindent\textbf{Antwort:} In the realm of combat, the ability to cast spells is generally not impeded by being grappled, unless other creatures are within 5 feet and disadvantage applies to the spellcasting ability check. However, making an attack or taking other actions, such as grappling or being grappled, may hinder the c\ldots

\noindent\textbf{Abgerufene Chunks:}
\begin{itemize}[noitemsep, topsep=2pt]
  \item[\textbf{\#1}] \texttt{[0.0167]} \textit{Monsters/# Monster Statistics} -- A monster's innate spells can't be swapped out with other spells. If a monster's\ldots
  \item[\textbf{\#2}] \texttt{[0.0164]} \textit{Spells/# Spellcasting} -- \# SPELLCASTING  Magic permeates fantasy gaming worlds and often appears in the f\ldots
  \item[\textbf{\#3}] \texttt{[0.0161]} \textit{Spells/# Spellcasting} -- Some characters and monsters have special abilities that let them cast spells wi\ldots
  \item[\textbf{\#4}] \texttt{[0.0159]} \textit{Gameplay/Combat} -- A grappled creature can use its action to escape. To do so, it must succeed on a\ldots
  \item[\textbf{\#5}] \texttt{[0.0156]} \textit{Gameplay/Combat} -- If either weapon has the thrown property, you can throw the weapon, instead of m\ldots
\end{itemize}
\vspace{4pt}
\hrule
\vspace{6pt}

\paragraph{Q12: If a creature has exhaustion level 3 and is also restrained, what is its movement speed?}
\noindent \textbf{Corr.:}~1.00/5\quad \textbf{Faith.:}~1.00\quad \textbf{Relev.:}~0.00\quad \textbf{Ctx.Rel.:}~0.62\\

\noindent\textbf{Antwort:} The context does not provide enough information to determine the movement speed of a creature with exhaustion level 3 and restrained. The rules regarding exhaustion only specify that a creature suffering from exhaustion has its speed halved at level 2. No information is given about how exhaustion le\ldots

\noindent\textbf{Abgerufene Chunks:}
\begin{itemize}[noitemsep, topsep=2pt]
  \item[\textbf{\#1}] \texttt{[0.0167]} \textit{Gamemastering/Conditions} -- A creature suffers the effect of its current level of exhaustion as well as all \ldots
  \item[\textbf{\#2}] \texttt{[0.0164]} \textit{Gamemastering/Conditions} -- Its weight increases by a factor of ten, and it ceases aging. - The creature is \ldots
  \item[\textbf{\#3}] \texttt{[0.0161]} \textit{Gameplay/Adventuring} -- For example, a creature with a Constitution of 14 can hold its breath for 3 minu\ldots
  \item[\textbf{\#4}] \texttt{[0.0159]} \textit{Gamemastering/Conditions} -- Also, being raised from the dead reduces a creature's exhaustion level by 1.  \#\#\ldots
  \item[\textbf{\#5}] \texttt{[0.0156]} \textit{Spells/Freedom of Movement} -- For the duration, the target's movement is unaffected by difficult terrain, and \ldots
\end{itemize}
\vspace{4pt}
\hrule
\vspace{6pt}

\paragraph{Q16: Can you use Action Surge to cast two leveled spells in one turn?}
\noindent \textbf{Corr.:}~2.00/5\quad \textbf{Faith.:}~0.00\quad \textbf{Relev.:}~0.00\quad \textbf{Ctx.Rel.:}~0.00\\

\noindent\textbf{Antwort:} I cannot find this in the provided rules about Action Surge or casting spells with multiple levels in one turn. However, a Wizard can cast multiple spells during one turn using their Action Surge, but each spell must have a different casting time and cannot be an evocation spell that affects other c\ldots

\noindent\textbf{Abgerufene Chunks:}
\begin{itemize}[noitemsep, topsep=2pt]
  \item[\textbf{\#1}] \texttt{[0.0167]} \textit{Classes/Wizard} -- When you cast an evocation spell that affects other creatures that you can see, \ldots
  \item[\textbf{\#2}] \texttt{[0.0164]} \textit{Spells/# Spellcasting} -- You can't concentrate on two spells at once. - **Taking damage.** Whenever you t\ldots
  \item[\textbf{\#3}] \texttt{[0.0161]} \textit{Classes/Fighter} -- On your turn, you can use a bonus action to regain hit points equal to 1d10 + yo\ldots
  \item[\textbf{\#4}] \texttt{[0.0159]} \textit{Spells/# Spellcasting} -- You can't cast another spell during the same turn, except for a cantrip with a c\ldots
  \item[\textbf{\#5}] \texttt{[0.0156]} \textit{Spells/Slow} -- On its turn, it can use either an action or a bonus action, not both. Regardless\ldots
\end{itemize}
\vspace{4pt}
\hrule
\vspace{6pt}
